\chapter{Plan de trabajo}
Se conserva el horizonte general del proyecto original: cierre del Seminario de Título en diciembre de 2026 y desarrollo y validación durante el primer semestre de 2027. La planificación se actualiza al 4 de octubre de 2026. Las fechas siguientes son estimaciones de trabajo, no confirmación de actividades realizadas.

\begin{table}[htbp]
\centering\small
\begin{tabularx}{\textwidth}{p{2.4cm}Xp{3.8cm}}
\toprule
Período & Actividad y objetivo & Entregable verificable \\
\midrule
Octubre 2026 & Depurar las 43 fichas y completar el protocolo bibliográfico (OE1). & Registro de búsquedas, exclusiones y metadatos. \\
Octubre--noviembre 2026 & Auditar datos, definir particiones y métricas (OE2). & Manifiesto de datos y protocolo reproducible. \\
Noviembre--diciembre 2026 & Implementar líneas base y piloto de medición (OE2--OE3). & Informe de avance con mediciones identificadas como piloto. \\
Enero--febrero 2027 & Desarrollar selección y comparación profesor--estudiante (OE4). & Módulos y resultados de ablación inicial. \\
Marzo--abril 2027 & Integrar adaptación y auditoría; evaluar caché si corresponde (OE4). & Prototipo y configuración congelada. \\
Abril--mayo 2027 & Ejecutar prueba final y análisis de calidad y recursos (OE5). & Resultados, intervalos y registro experimental. \\
Junio 2027 & Redactar discusión, conclusiones y preparar defensa. & Memoria y material de presentación. \\
\bottomrule
\end{tabularx}
\caption{Plan actualizado por entregables. La ejecución de etapas posteriores depende de la auditoría de datos y del piloto.}
\end{table}

Al cierre de esta reformulación se dispone del corpus organizado y del diseño experimental escrito. No se considera completada la revisión sistemática ni se atribuyen avances de implementación sin evidencia de código y mediciones. El primer hito técnico será obtener una línea base reproducible; solo después se incorporarán los niveles de optimización.

Si el piloto muestra que TreeSHAP completo satisface el presupuesto disponible, la investigación comparará ese resultado con las variantes adaptativas. Si la destilación no aporta una ventaja, se conservará su evaluación como resultado y se priorizará la configuración más sencilla que cumpla las restricciones. La planificación admite estas decisiones sin cambiar la pregunta central.
